\documentclass{article}
\usepackage[utf8]{inputenc}
\usepackage[T1]{fontenc}
\usepackage{geometry}
\geometry{
    a4paper,
    left=2.5cm,
    right=2.5cm,
    top=2.5cm,
    bottom=2.5cm
}
\usepackage{graphicx}
\usepackage{float}
\usepackage{hyperref}
\usepackage[english]{babel}
\usepackage{amsmath}
\usepackage{amssymb}
\usepackage{array}
\usepackage{booktabs}
\usepackage{caption}
\usepackage{subcaption}
\usepackage{enumitem}
\usepackage{fancyhdr}
\usepackage{titlesec}
\usepackage{url}
\usepackage{textcomp}
\usepackage{xcolor}
\usepackage{listings}
\usepackage{multirow}
\usepackage{longtable}
\usepackage{tabularx}
\usepackage{times}

% Set up page style
\pagestyle{fancy}
\fancyhf{}
\rfoot{\thepage}
\renewcommand{\headrulewidth}{0pt}

% Title formatting
\titleformat{\chapter}[display]{\normalfont\bfseries}{\chaptertitle\thechapter}{0pt}{\Huge}
\titlespacing*{\chapter}{0pt}{50pt}{40pt}

% Document info
\title{Template Document}
\author{Academic Support Agent}
\date{\today}

\begin{document}
\maketitle

%\pretextual%% Elementos pr -textuais

\begin{paginadetitulo}%% P gina de t tulo

\begin{ambienteresumo}%% Resumo
Voc  deve estruturar o resumo de modo a descrever o seu trabalho com simplicidade e objetividade. Tenha sempre em mente que o resumo   a apresenta o do seu trabalho a leitores que provavelmente n o est o familiarizados com alguns termos e express es que podem ser comuns em sua  rea de estudo. O texto do resumo deve ser apresentado em um  nico par grafo, sem recuo na primeira linha. A fonte usada deve ser Arial, com 10 pontos de tamanho e alinhamento justificado. O resumo deve ter entre 150 e 200 palavras: por exemplo, este texto possui exatamente 161 palavras. Na medida do poss vel, evite usar f rmulas, s mbolos matem ticos e refer ncias bibliogr ficas em seu resumo. Se ambos os resumos ultrapassarem a p gina inicial do documento, n o h  problema. Preferencialmente, procure apresentar o objetivo, o m todo de pesquisa, os principais resultados e as conclus es poss veis a partir do seu trabalho. Um bom resumo informa de maneira clara e objetiva o que ser  apresentado no texto.
\palavraschave{palavra um;\ palavra dois;\ palavra tr s. (em  ordem alfab tica).}%% Palavras-chave
\end{ambienteresumo}

\begin{ambienteresumo}[Abstract]%% Abstract
\begin{otherlanguage*}{english}%% Idioma do abstract
The abstract of your work, in addition to being part of this document, must be inserted in the Abstract field in the third step of submission on the portal https://eventos.utfpr.edu.br/. You should structure the abstract to describe your work simply and objectively. Always keep in mind that the abstract is the presentation of your work to readers who are probably unfamiliar with some terms and expressions. The text of the abstract must be presented in a single paragraph. The font used must be Arial, with 10-point size and justified alignment. The number of words may be modified due to the translation, for example, this text has exactly 167 words. As far as possible, avoid using formulas, mathematical symbols, and references in your abstract. It is okay if the two abstracts exceed the front page of the document. Preferably, try to present the objective, research method, main results and possible conclusions from your work. A good summary clearly and objectively informs what will be presented in the text.
%Follow the same instructions presented on section \textit{RESUMO}. Translate here your \textit{RESUMO} to English. 
\palavraschave[KEYWORDS]{keyword one; keyword two; keyword three. (Keep the order of words in Portuguese).}%% Keywords

\blfootnote{\hfill \color{gray} \sffamily \large{seisicite.com.br \hspace{1cm} \thepage}}

\end{otherlanguage*}
\end{ambienteresumo}
\end{paginadetitulo}
\textual%% Elementos textuais

\section{T TULO DE SE O DE N VEL 1}\label{sec:intro}

A estrutura do trabalho pode ser dividida em se es   sua escolha, por exemplo: Introdu o, Materiais e M todos, Resultados e Discuss es, Conclus o, Agradecimentos e Refer ncias. Os itens obrigat rios s o informa es sobre financiamento, conflito de interesse e refer ncias.

Os T tulos das se es devem ser posicionados   esquerda e sem ponto final, com fonte Arial em caixa alta e negrito (como mostrado acima). Os t tulos das se es de n vel 2 seguem a mesma formata o, exceto pelo n o uso do negrito (como mostrado abaixo).

\subsection{T TULO DE SE O DE N VEL 2}

O texto completo deve ser redigido em portugu s e conter entre quatro e seis p ginas de tamanho A4, incluindo nesta contagem as p ginas com t tulo, autores, resumo, abstract e refer ncias.

O corpo do texto deve apresentar fonte Arial de tamanho 12, alinhamento justificado, sendo a primeira linha de cada par grafo deslocada 1,25 cm.

Deve ser aplicado espa amento simples entre as linhas em todo o trabalho e sem espa amento entre par grafos. As p ginas devem ser numeradas na parte inferior   direita conforme o modelo deste documento.

As cita es devem obedecer ao sistema autor-data e estar de acordo com a norma NBR 10520 da Associa o Brasileira de Normas T cnicas (ABNT). Cita es diretas de at  tr s linhas devem estar entre aspas e no fluxo do texto. Cita es com mais de tr s linhas devem aparecer a partir de um deslocamento de 4 cm (a partir da margem esquerda) sem aspas.

As refer ncias utilizadas ao longo do trabalho devem ser listadas no final do trabalho de acordo com as instru es na se o Refer ncias. Sugerimos o uso de programas de gest o de refer ncias bibliogr ficas como o Zotero para garantir que as regras da ABNT sejam seguidas.

A reda o do texto pode ser feita em qualquer editor de textos como Google Docs, MS Word, LO Writer ou Overleaf (no caso de utilizar template em \LaTeX). Por m, o arquivo deve ser convertido em PDF para a submiss o, para evitar problemas de formata o por incompatibilidade de vers es de programas. N o ser o aceitos para avalia o trabalhos em outro formato.

Qualquer que seja o tipo de ilustra o inserida no trabalho (fluxograma, gr fico, quadro, figura, imagem, tabela, entre outros), sua identifica o deve aparecer na parte superior, em negrito, precedida da palavra designativa, seguida de seu n mero de ordem de ocorr ncia no texto, em algarismos ar bicos, travess o e do respectivo t tulo e com fonte de tamanho 10.

Ap s a ilustra o, na parte inferior, indicar a fonte consultada (mesmo sendo produ o do pr prio autor), legenda, notas e outras informa es necess rias   sua compreens o (se houver).

A ilustra o deve ser citada no texto e inserida o mais pr ximo poss vel do trecho a que se refere. Ver, por exemplo, a Figura \ref{fig:grafico}.

\begin{figure}[H]
\centering%
\caption{Exemplo de figura}%
\label{fig:grafico}
\includegraphics[width = 0.525\columnwidth]{./Figuras/grafico}
\fonte{Indicar a fonte consultada (mesmo sendo produ o do pr prio autor).}
\end{figure}

%Figuras podem ser inseridas neste documento usando o ambiente \LaTeX\ \enquote{figure}, conforme exemplo no arquivo fonte deste modelo.

%\subsubsection{Tabelas e quadros}\label{sssec:tabs}

Tabelas e quadros devem estar centralizados e conter apenas dados imprescind veis, evitando-se que sejam muito extensos, n o repetindo dados j  inseridos no texto, ou vice-versa. O formato pode ser observado na Tabela \ref{tab:resul}.

\begin{table*}[!htb]
\centering%
\normalsize
\caption{Exemplo de tabela.}%
\label{tab:resul}
\begin{tabular}{l | c }
\toprule
Idade   & Percentual \%    \\
\midrule
At  20 anos & 0\\
Entre 21 e 30 anos & 10\\
Entre 31 e 40 anos & 20\\
Entre 41 e 50 anos & 30\\
Acima de 51 anos & 40\\
\bottomrule
\addlinespace
\end{tabular}
\fonte{Beltrano (2021).}
\end{table*}

%Para o uso de quadros, siga a estrutura demonstrada no Quadro \ref{quad:perfil}. Caso os dados sejam in ditos, provenientes de uma pesquisa dos pr prios autores, use a fonte Autoria pr pria, seguida do ano que o quadro foi elaborado com base em \cite{IBGE2017}.

No caso de quadros, deve ser seguida a estrutura demonstrada no Quadro \ref{quad:secoes}. Caso os dados sejam in ditos e provenientes de uma pesquisa realizada pelos pr prios autores do trabalho, essa especifica o deve constar na fonte com o ano da pesquisa de campo. Nesse caso a fonte deve ser: Elaborado pelos autores (2021). Caso o material seja elaborado a partir de dados secund rios, informar: Elaborado pelos autores com base em \cite{IBGE2017}.

%Deixe uma linha em branco entre o quadro e o texto.

\begin{quadro}
\centering%
\textbf{Exemplo de quadro.}\\%
\begin{tabular}{| c | c | c |}
\hline
\textbf{Se o}    & \textbf{Tipografia}    & \textbf{Exemplo}     \\
\hline
Se es prim rias & Letras mai sculas em negrito    & \textbf{SE O PRIM RIA}   \\
\hline
Se es secund rias & Letras mai sculas sem negrito    & SE O SECUND RIA   \\
\hline
\end{tabular}
\vspace{.15cm}
\fonte{Elaborado por autores (2022).}
\label{quad:secoes}
\end{quadro}

\subsubsection{Material Suplementar}
As equa es matem ticas devem aparecer a partir de um deslocamento de 0,6 cm (a partir da margem esquerda). N meros ar bicos devem ser usados em equa es, inseridos entre par nteses, como ilustrado na \cref{eq:Uxy}.
As equa es devem ser referenciadas por \cref{eq:Uxy} no meio ou no final da frase, ou por \Cref{eq:Uxy} quando usada no in cio de uma senten a.

\begin{equation}\label{eq:Uxy}
u = \beta \sin(\pi x)\frac{(e^{2x}-1)(e^y-1)}{(e^2-1)(e-1)}
\end{equation}

\subsubsection{Agradecimentos}

Esta se o   \underline{obrigat ria \textbf{na vers o com autores}}, N O INCLUIR na vers o sem autores. Ela serve para agradecer as institui es e pessoas que contribu ram para a realiza o do trabalho, especialmente  rg os fomento que viabilizaram recursos no formato de bolsa e aux lio financeiro. 

\subsubsection{Disponibilidade de C digo}

Esta se o   opcional. Ela serve para declarar se o c digo desenvolvido est  dispon vel para terceiros. Caso esteja, explicitar onde pode ser encontrado. Caso contr rio, explicitar a raz o para a indisponibilidade. Remova caso nenhum c digo tenha sido desenvolvido. 

\subsubsection{Conflito de interesse}

Esta se o   \underline{obrigat ria}. Ela serve para os autores declararem se possuem ou n o algum conflito de interesse. N o havendo, sugerimos declarar  N o h  conflito de interesse .

\begin{center}
    \textbf{$\ast$ Instru es para se o de refer ncias}
\end{center}

%\raggedright{\textbf{Instru es para se o de refer ncias}}

A se o de Refer ncias (abaixo)   \underline{obrigat ria}. Todas as obras citadas no texto devem ser apresentadas no final do documento. As refer ncias devem ser baseadas na norma \href{http://portal.utfpr.edu.br/alunos/normas-academicas/normas-para-trabalhos-academicos}{NBR 6023/2002}, norma da ABNT atualmente em vigor. Como sugest o utilize o Sistema de Gest o de Normas e Documentos Regulat rios (GEDWEB) da UTFPR (\url{https://webapp.utfpr.edu.br/bibservices/gedWeb}) e um software de gest o de refer ncias bibliogr ficas como o Zotero.

\textbf{Importante:} as refer ncias ser o alinhadas   esquerda do texto (n o justificado), em ordem alfab tica, espa amento simples entre linhas, com espa amento de 12 pontos antes e de 6 pontos depois. A ABNT recomenda seguir uma uniformidade. Portanto, defina se voc  utilizar  o nome por extenso dos autores ou se ir  abrevi -los. Ao definir, siga o mesmo padr o para todas as refer ncias. 
Exemplo ilustrativo na se o de Refer ncias (abaixo), onde se utiliza  nome por extenso dos autores conforme as seguintes cita es: 
\cite{volpato2016autoria},
\cite{schneider2020efeitos} e
\cite{MEDEIROS_2021}.

% Refer ncias
\printbibliography[title={\raggedright{Refer ncias}}]

%% Fim do documento
\end{document}